% ============================================
% Section: Introduction (MobiCom-focused, wireless-centric)
% ============================================

Cooperative autonomous systems---from multi-robot construction teams to
UAV--UGV disaster-response squads---require a shared, up-to-date 3D model of
their environment to plan, navigate, and coordinate. Recent advances in 3D
Gaussian Splatting (3DGS)~\cite{kerbl20233dgs} have made it possible to
construct photo-realistic maps in real time from a single mobile
platform~\cite{lang2026gaussian2,lang2025gaussian,xie2025gslivm,park2026livegs},
and multi-agent Gaussian SLAM has emerged as an active research
frontier~\cite{deambrogi2026cgsslam_mob,thomas2025grandslam_mob,deng2025mneslam}.

Yet a critical gap remains: \textbf{no existing system addresses the wireless
communication problem that arises when multiple nodes attempt to share
photo-realistic Gaussian maps over real, bandwidth-limited, lossy wireless
channels.} A typical submap from a single node traveling 5~m contains
50,000--100,000 3D Gaussians. Even with aggressive compression, a single submap
consumes 2--5~MB---and a cluster of 8 nodes publishing every 3~s generates
5--13~Mbps of aggregate uplink traffic on the cluster head. Over Wi-Fi or 5G,
this is manageable; over degraded links (congested Wi-Fi channels, 5G edge-cell
handover, ad-hoc mesh), it is not. Without careful protocol design, submap
delivery fails, registrations stall, and the global map diverges.

This paper presents \textbf{Co-LIC2}, a cooperative multi-node
LiDAR-inertial-visual Gaussian SLAM system whose central contribution is a
purpose-built wireless communication architecture that jointly addresses
\emph{data representation}, \emph{protocol design}, and \emph{network topology}
for photo-realistic map sharing. Co-LIC2 extends the single-platform
Gaussian-LIC2 framework~\cite{lang2026gaussian2} along four axes:

\begin{enumerate}
  \item \textbf{A QoS-aware multi-stream wireless protocol}
  (Section~\ref{sec:protocol}) that classifies traffic into four streams
  (real-time pose, bulk submap, reference keyframe, control) and assigns
  differentiated 802.11e EDCA parameters, application-layer FEC, and
  deadline-aware scheduling. This is the \emph{core contribution} to the mobile
  computing community: a protocol specification that can be implemented on
  commodity Wi-Fi~6/6E radios without modifying the MAC layer.

  \item \textbf{A submap-aware adaptive compression scheduler}
  (Section~\ref{sec:compression}) that jointly optimizes SH-degree truncation,
  spatial quantization step, and zstd entropy-coding level under a hard
  bandwidth cap, using a new perceptually motivated Number-of-Gaussians-
  Per-View (NPV) metric as the objective function, rather than naive byte-budget
  heuristics.

  \item \textbf{A hierarchical hybrid network architecture}
  (Section~\ref{sec:network}) with (a) a stable cluster-head election protocol
  that accounts for compute capacity, link quality, and
  sensor richness; (b) multi-hop submap routing with store-and-forward on
  partition; and (c) three degradation modes (full cooperation $\to$ intra-cluster
  P2P $\to$ autonomous single-node) that preserve map consistency under
  intermittent connectivity.

  \item \textbf{A cooperative pose-graph SLAM back-end}
  (Section~\ref{sec:mapping}) that registers, merges, and globally optimizes
  ownership-tagged Gaussian submaps from heterogeneous nodes with confidence-
  weighted conflict resolution.
\end{enumerate}

We provide a complete protocol specification with finite-state machines, a
submap binary format, analytical latency and scalability bounds, and a
detailed failure-mode recovery analysis. Co-LIC2's design is validated against
realistic deployment scenarios: 2--32 heterogeneous nodes, per-node uplink
$\geq$3.2~Mbps over Wi-Fi~6, per-cluster submap merge latency $\leq$2~s,
and graceful degradation to single-node Gaussian-LIC2 under total partition.
An implementation blueprint and evaluation plan on physical multi-rig hardware
are provided.
